\documentclass[12pt,a4paper]{article}
\usepackage[utf8]{inputenc}
\usepackage[spanish]{babel}
\usepackage{geometry}
\usepackage{listings}
\usepackage{xcolor}
\usepackage{hyperref}
\usepackage{graphicx}
\usepackage{fancyhdr}
\usepackage{tocloft}

\geometry{margin=2.5cm}

% Configuración de listings para código
\lstset{
    basicstyle=\ttfamily\small,
    backgroundcolor=\color{gray!10},
    frame=single,
    rulecolor=\color{black!30},
    breaklines=true,
    commentstyle=\color{green!60!black},
    keywordstyle=\color{blue},
    stringstyle=\color{red},
    showstringspaces=false,
    captionpos=b
}

% Configuración de headers y footers
\pagestyle{fancy}
\fancyhf{}
\fancyhead[L]{Guía Git para Desarrolladores Java}
\fancyhead[R]{\thepage}
\fancyfoot[C]{De SVN a Git - Transición Completa}

\title{{\Huge \textbf{Git para Desarrolladores Java}} \\ 
       {\Large Transición desde SVN con Ejemplos Prácticos}}
\author{Guía Completa para Desarrolladores con 1 año de experiencia}
\date{\today}

\begin{document}

\maketitle

\tableofcontents
\newpage

\section{Introducción: ¿Por qué cambiar de SVN a Git?}

Como desarrollador Java con experiencia en SVN, ya tienes una base sólida en control de versiones. Git te ofrecerá nuevas capacidades que transformarán tu flujo de trabajo:

\subsection{Principales diferencias}
\begin{itemize}
    \item \textbf{Distribución}: Cada clon es un repositorio completo
    \item \textbf{Ramas}: Extremadamente ligeras y rápidas
    \item \textbf{Offline}: Puedes trabajar sin conexión al servidor
    \item \textbf{Staging Area}: Área intermedia para preparar commits
    \item \textbf{Historia}: Historial completo en tu máquina local
\end{itemize}

\subsection{Comparación SVN vs Git}

\begin{center}
\begin{tabular}{|l|l|l|}
\hline
\textbf{Aspecto} & \textbf{SVN} & \textbf{Git} \\
\hline
Arquitectura & Centralizada & Distribuida \\
Ramas & Costosas (carpetas) & Ligeras (punteros) \\
Commits & Directos al servidor & Locales + Push \\
Trabajo offline & Limitado & Completo \\
Velocidad & Depende de red & Muy rápido \\
\hline
\end{tabular}
\end{center}

\section{Conceptos Fundamentales}

\subsection{Los Tres Estados de Git}
\begin{enumerate}
    \item \textbf{Working Directory}: Archivos en tu directorio de trabajo
    \item \textbf{Staging Area}: Archivos preparados para commit
    \item \textbf{Repository}: Historial de commits confirmados
\end{enumerate}

\begin{lstlisting}[caption=Flujo básico de estados]
# Working Directory -> Staging Area
git add archivo.java

# Staging Area -> Repository
git commit -m "Mensaje del commit"

# Repository -> Servidor remoto
git push origin main
\end{lstlisting}

\subsection{Ramas en Git}
A diferencia de SVN donde las ramas son carpetas costosas, en Git son simplemente punteros ligeros.

\begin{lstlisting}[caption=Manejo básico de ramas]
# Crear nueva rama
git branch feature/nueva-funcionalidad

# Cambiar a la rama
git checkout feature/nueva-funcionalidad

# Crear y cambiar en un comando
git checkout -b feature/otra-funcionalidad

# Listar ramas
git branch

# Fusionar rama
git checkout main
git merge feature/nueva-funcionalidad
\end{lstlisting}

\section{Instalación y Configuración Inicial}

\subsection{Configuración básica}
\begin{lstlisting}[caption=Configuración inicial de Git]
# Configurar identidad
git config --global user.name "Tu Nombre"
git config --global user.email "tu.email@empresa.com"

# Configurar editor
git config --global core.editor "code --wait"

# Ver configuración
git config --list
\end{lstlisting}

\subsection{Configuración específica para Java}
\begin{lstlisting}[caption=Configuraciones útiles para desarrollo Java]
# Ignorar archivos compilados por defecto
git config --global core.excludesfile ~/.gitignore_global

# Configurar diff para archivos Java
git config --global diff.tool vscode

# Configurar merge tool
git config --global merge.tool vscode
\end{lstlisting}

\section{Comandos Esenciales con Ejemplos}

\subsection{Inicialización y Clonado}
\begin{lstlisting}[caption=Crear y clonar repositorios]
# Inicializar repositorio nuevo
git init mi-proyecto-java

# Clonar repositorio existente
git clone https://github.com/usuario/proyecto-java.git

# Clonar con nombre específico
git clone https://github.com/usuario/proyecto.git mi-proyecto
\end{lstlisting}

\subsection{Seguimiento de Archivos}
\begin{lstlisting}[caption=Gestión de archivos]
# Ver estado de archivos
git status

# Agregar archivos específicos
git add src/main/java/com/empresa/MiClase.java

# Agregar todos los archivos Java
git add "*.java"

# Agregar todo
git add .

# Ver diferencias
git diff
git diff --staged
\end{lstlisting}

\subsection{Commits}
\begin{lstlisting}[caption=Creación de commits]
# Commit básico
git commit -m "Implementar validación de usuario"

# Commit con descripción detallada
git commit -m "Implementar validación de usuario

- Agregar clase ValidadorUsuario
- Implementar validación de email
- Agregar tests unitarios
- Actualizar documentación"

# Commit que incluye modificaciones
git commit -a -m "Corregir bug en login"
\end{lstlisting}

\section{Flujos de Trabajo Comunes}

\subsection{Git Flow para Equipos Java}
\begin{lstlisting}[caption=Flujo típico en proyecto Java]
# 1. Actualizar rama principal
git checkout main
git pull origin main

# 2. Crear rama para nueva funcionalidad
git checkout -b feature/JIRA-123-validacion-usuarios

# 3. Desarrollar y hacer commits
# ... modificar archivos ...
git add .
git commit -m "JIRA-123: Implementar validación básica"

# 4. Subir rama al servidor
git push -u origin feature/JIRA-123-validacion-usuarios

# 5. Crear Pull Request en la plataforma
# 6. Después de aprobación, fusionar
git checkout main
git pull origin main
git branch -d feature/JIRA-123-validacion-usuarios
\end{lstlisting}

\subsection{Hotfixes}
\begin{lstlisting}[caption=Manejo de correcciones urgentes]
# Crear rama de hotfix desde main
git checkout main
git pull origin main
git checkout -b hotfix/correccion-critica-login

# Hacer la corrección
# ... modificar archivos ...
git add .
git commit -m "Hotfix: Corregir NPE en login"

# Subir y fusionar urgentemente
git push -u origin hotfix/correccion-critica-login
# Crear PR urgente para revisión rápida
\end{lstlisting}

\section{Archivos .gitignore para Proyectos Java}

\subsection{Contenido típico de .gitignore}
\begin{lstlisting}[caption=.gitignore para proyectos Java]
# Archivos compilados
*.class
*.jar
*.war
*.ear

# IDEs
.idea/
.vscode/
*.iml
*.iws
*.ipr

# Eclipse
.metadata
.classpath
.project
.settings/

# Maven
target/
pom.xml.tag
pom.xml.releaseBackup
pom.xml.versionsBackup

# Gradle
.gradle/
build/

# Logs
*.log
logs/

# OS
.DS_Store
Thumbs.db

# Spring Boot
application-local.properties
application-dev.properties
\end{lstlisting}

\section{Resolución de Conflictos}

\subsection{Conflictos de Merge}
\begin{lstlisting}[caption=Resolver conflictos paso a paso]
# 1. Intentar merge
git merge feature/nueva-funcionalidad

# Si hay conflictos, Git mostrará:
# Auto-merging src/main/java/com/empresa/Usuario.java
# CONFLICT (content): Merge conflict in Usuario.java

# 2. Ver archivos con conflictos
git status

# 3. Editar archivos manualmente
# Buscar marcadores <<<<<<< ======= >>>>>>>

# 4. Marcar como resueltos
git add src/main/java/com/empresa/Usuario.java

# 5. Completar merge
git commit -m "Resolver conflicto en clase Usuario"
\end{lstlisting}

\section{Comandos Avanzados Útiles}

\subsection{Historial y Navegación}
\begin{lstlisting}[caption=Explorar historial]
# Ver historial
git log --oneline --graph --all

# Ver cambios en archivo específico
git log -p src/main/java/com/empresa/Usuario.java

# Buscar en historial
git log --grep="JIRA-123"

# Ver quién modificó cada línea
git blame src/main/java/com/empresa/Usuario.java
\end{lstlisting}

\subsection{Deshacer Cambios}
\begin{lstlisting}[caption=Deshacer diferentes tipos de cambios]
# Descartar cambios no confirmados
git checkout -- archivo.java

# Deshacer staging
git reset HEAD archivo.java

# Deshacer último commit (mantener cambios)
git reset --soft HEAD~1

# Deshacer último commit (descartar cambios)
git reset --hard HEAD~1

# Revertir commit público
git revert abc123def
\end{lstlisting}

\section{Integración con IDEs Java}

\subsection{IntelliJ IDEA}
\begin{itemize}
    \item \textbf{VCS Menu}: Acceso completo a operaciones Git
    \item \textbf{Git Tool Window}: Vista del historial y ramas
    \item \textbf{Commit Tool Window}: Staging y commits visuales
    \item \textbf{Merge Conflicts}: Herramienta visual de resolución
\end{itemize}

\subsection{Eclipse}
\begin{itemize}
    \item \textbf{EGit Plugin}: Integración completa con Git
    \item \textbf{Git Repositories View}: Gestión de repositorios
    \item \textbf{Git Staging View}: Control del staging area
    \item \textbf{History View}: Exploración del historial
\end{itemize}

\section{Mejores Prácticas para Equipos Java}

\subsection{Convenciones de Commits}
\begin{lstlisting}[caption=Formato de mensajes de commit]
# Formato: tipo(scope): descripción

feat(auth): implementar validación JWT
fix(login): corregir NPE en validación
docs(api): actualizar documentación REST
test(user): agregar tests de integración
refactor(db): optimizar consultas de usuario
\end{lstlisting}

\subsection{Estructura de Ramas}
\begin{itemize}
    \item \textbf{main/master}: Código de producción
    \item \textbf{develop}: Integración de funcionalidades
    \item \textbf{feature/*}: Nuevas funcionalidades
    \item \textbf{hotfix/*}: Correcciones urgentes
    \item \textbf{release/*}: Preparación de versiones
\end{itemize}

\section{Migración desde SVN}

\subsection{Preparación}
\begin{lstlisting}[caption=Pasos para migrar desde SVN]
# 1. Crear archivo de mapeo de usuarios
echo "usuario_svn = Nombre Completo <email@empresa.com>" > users.txt

# 2. Clonar repositorio SVN
git svn clone --authors-file=users.txt \
  http://svn.empresa.com/repo/proyecto \
  --trunk=trunk --branches=branches --tags=tags

# 3. Convertir ramas y tags
cd proyecto
git for-each-ref refs/remotes/origin | cut -d/ -f3- | \
  grep -v '^tags/' | while read branchname; do
    git branch "$branchname" "refs/remotes/origin/$branchname"
    git branch -r -d "origin/$branchname"
done

# 4. Crear repositorio Git remoto
git remote add origin https://github.com/empresa/proyecto.git
git push -u origin main
\end{lstlisting}

\section{Herramientas Complementarias}

\subsection{GUI Tools}
\begin{itemize}
    \item \textbf{SourceTree}: GUI gratuita multiplataforma
    \item \textbf{GitKraken}: GUI moderna con funciones avanzadas
    \item \textbf{GitHub Desktop}: Simple para repositorios GitHub
    \item \textbf{TortoiseGit}: Integración con Windows Explorer
\end{itemize}

\subsection{Plataformas de Hosting}
\begin{itemize}
    \item \textbf{GitHub}: Más popular, gran ecosistema
    \item \textbf{GitLab}: CI/CD integrado, versión self-hosted
    \item \textbf{Bitbucket}: Integración con Atlassian (Jira)
    \item \textbf{Azure DevOps}: Integración con Microsoft Stack
\end{itemize}

\section{Troubleshooting Común}

\subsection{Problemas Frecuentes}
\begin{lstlisting}[caption=Soluciones a problemas comunes]
# Problema: "fatal: refusing to merge unrelated histories"
git pull origin main --allow-unrelated-histories

# Problema: Archivos grandes accidentalmente agregados
git filter-branch --force --index-filter \
  'git rm --cached --ignore-unmatch archivo-grande.jar' \
  --prune-empty --tag-name-filter cat -- --all

# Problema: Cambiar mensaje del último commit
git commit --amend -m "Nuevo mensaje"

# Problema: Olvidé agregar archivo al último commit
git add archivo-olvidado.java
git commit --amend --no-edit
\end{lstlisting}

\section{Ejercicios Prácticos}

\subsection{Ejercicio 1: Flujo Básico}
\begin{enumerate}
    \item Crear repositorio para proyecto Java
    \item Configurar .gitignore apropiado
    \item Crear estructura Maven/Gradle básica
    \item Hacer commits incrementales
    \item Crear y fusionar rama de funcionalidad
\end{enumerate}

\subsection{Ejercicio 2: Colaboración}
\begin{enumerate}
    \item Simular conflicto de merge
    \item Practicar resolución manual
    \item Usar herramientas visuales del IDE
    \item Crear Pull Request
    \item Revisar código de compañero
\end{enumerate}

\section{Referencias y Recursos}

\subsection{Documentación Oficial}
\begin{itemize}
    \item Pro Git Book: \url{https://git-scm.com/book}
    \item Git Reference: \url{https://git-scm.com/docs}
    \item GitHub Guides: \url{https://guides.github.com/}
\end{itemize}

\subsection{Cursos y Tutoriales}
\begin{itemize}
    \item Atlassian Git Tutorial: \url{https://www.atlassian.com/git/tutorials}
    \item Learn Git Branching: \url{https://learngitbranching.js.org/}
    \item GitHub Learning Lab: \url{https://lab.github.com/}
\end{itemize}

\section{Conclusiones}

La transición de SVN a Git puede parecer compleja inicialmente, pero las ventajas son significativas:

\begin{itemize}
    \item \textbf{Velocidad}: Operaciones locales instantáneas
    \item \textbf{Flexibilidad}: Ramas ligeras y múltiples flujos de trabajo
    \item \textbf{Distribución}: Cada desarrollador tiene historial completo
    \item \textbf{Integración}: Mejor soporte en IDEs modernos
    \item \textbf{Ecosistema}: Herramientas y servicios abundantes
\end{itemize}

Con práctica constante y los ejemplos de esta guía, dominarás Git rápidamente y mejorarás significativamente tu productividad como desarrollador Java.

\section{Fundamentos Profundos de Git}

\subsection{Arquitectura Interna de Git}
Git no es solo un sistema de control de versiones, es un sistema de archivos dirigido por contenido con una interfaz VCS.

\subsubsection{El Modelo de Datos}
\begin{lstlisting}[caption=Estructura interna de objetos Git]
# Git almacena 4 tipos de objetos:

# 1. BLOB (Binary Large Object) - Contenido de archivos
# SHA-1: da39a3ee5e6b4b0d3255bfef95601890afd80709 (archivo vacío)

# 2. TREE - Estructura de directorios  
# Contiene referencias a blobs y otros trees

# 3. COMMIT - Instantánea del proyecto
# Referencia a tree + metadatos (autor, fecha, mensaje)

# 4. TAG - Referencia a commit con metadata adicional
\end{lstlisting}

\subsubsection{Hash SHA-1 y Integridad}
\begin{lstlisting}[caption=Verificación de integridad con SHA-1]
# Cada objeto Git tiene un hash SHA-1 único
# Basado en el contenido, garantiza integridad

# Ejemplo de hash de commit
git log --oneline
# a1b2c3d Mensaje del commit

# Ver el objeto completo
git cat-file -p a1b2c3d
# tree 4b825dc642cb6eb9a060e54bf8d69288fbee4904
# author Juan Pérez <juan@empresa.com> 1700000000 +0200
# committer Juan Pérez <juan@empresa.com> 1700000000 +0200
# 
# Mensaje del commit
\end{lstlisting}

\subsection{El Área de Preparación (Staging Area) en Detalle}
\begin{lstlisting}[caption=Manipulación avanzada del staging area]
# Ver contenido del índice
git ls-files --stage

# Añadir partes específicas de un archivo
git add -p archivo.java

# Añadir líneas específicas interactivamente
git add -i

# Resetear archivo del staging área
git reset HEAD archivo.java

# Ver diferencias entre working dir, staging y repository
git diff                    # working vs staging
git diff --staged          # staging vs repository  
git diff HEAD              # working vs repository
\end{lstlisting}

\section{Comandos Avanzados y Técnicas Expertas}

\subsection{Manipulación Avanzada del Historial}

\subsubsection{Git Rebase - Reescribir Historia}
\begin{lstlisting}[caption=Rebase interactivo y avanzado]
# Rebase interactivo para limpiar commits
git rebase -i HEAD~3

# Opciones disponibles en rebase interactivo:
# pick   = usar commit tal como está
# reword = usar commit pero editar mensaje
# edit   = usar commit pero parar para enmendar
# squash = usar commit pero fusionar con el anterior
# fixup  = como squash pero descartar mensaje
# drop   = remover commit

# Rebase con estrategia específica
git rebase -X theirs feature-branch

# Rebase preservando merges
git rebase --rebase-merges main

# Continuar rebase después de resolver conflictos
git rebase --continue

# Abortar rebase
git rebase --abort
\end{lstlisting}

\subsubsection{Git Cherry-pick Avanzado}
\begin{lstlisting}[caption=Cherry-pick con opciones avanzadas]
# Cherry-pick múltiples commits
git cherry-pick commit1 commit2 commit3

# Cherry-pick rango de commits
git cherry-pick start-commit..end-commit

# Cherry-pick sin hacer commit automático
git cherry-pick -n commit-hash

# Cherry-pick editando el mensaje
git cherry-pick -e commit-hash

# Cherry-pick firmando el commit
git cherry-pick -S commit-hash
\end{lstlisting}

\subsection{Búsqueda y Depuración Avanzada}

\subsubsection{Git Bisect - Búsqueda Binaria de Bugs}
\begin{lstlisting}[caption=Encontrar bugs con bisect]
# Iniciar bisect
git bisect start

# Marcar commit actual como malo
git bisect bad

# Marcar commit conocido bueno
git bisect good v1.0.0

# Git automáticamente selecciona commit del medio
# Probar y marcar como bueno o malo
git bisect good  # si funciona
git bisect bad   # si tiene el bug

# Automatizar con script
git bisect run test-script.sh

# Finalizar bisect
git bisect reset
\end{lstlisting}

\subsubsection{Git Grep Avanzado}
\begin{lstlisting}[caption=Búsqueda avanzada en historial]
# Buscar en todo el repositorio
git grep -n "patrón"

# Buscar en commits específicos
git grep "patrón" commit-hash

# Buscar con regex
git grep -E "patrón|otro-patrón"

# Buscar considerando contexto
git grep -A 3 -B 3 "patrón"

# Buscar en archivos específicos
git grep "patrón" -- "*.java"

# Buscar ignorando case
git grep -i "patrón"

# Contar ocurrencias
git grep -c "patrón"
\end{lstlisting}

\subsection{Configuración Avanzada de Git}

\subsubsection{Configuraciones de Performance}
\begin{lstlisting}[caption=Optimizaciones de rendimiento]
# Configurar compresión
git config --global pack.compression 9

# Configurar ventana de delta
git config --global pack.window 250

# Habilitar preload index
git config --global core.preloadindex true

# Configurar parallel index
git config --global core.untrackedCache true

# Configurar fsmonitor (para repositorios grandes)
git config --global core.fsmonitor true

# Configurar split-index
git config --global core.splitIndex true
\end{lstlisting}

\subsubsection{Hooks Avanzados}
\begin{lstlisting}[caption=Scripts de automatización con hooks]
# Pre-commit hook para validar código Java
#!/bin/bash
# .git/hooks/pre-commit

# Ejecutar checkstyle
mvn checkstyle:check
if [ $? -ne 0 ]; then
    echo "Checkstyle falló. Corrige los errores antes del commit."
    exit 1
fi

# Ejecutar tests unitarios
mvn test
if [ $? -ne 0 ]; then
    echo "Tests fallan. Corrige antes del commit."
    exit 1
fi

# Pre-push hook para validar branch
#!/bin/bash
# .git/hooks/pre-push

protected_branch='main'
current_branch=$(git symbolic-ref HEAD | sed -e 's,.*/\(.*\),\1,')

if [ $current_branch = $protected_branch ]; then
    echo "Push directo a $protected_branch no permitido."
    exit 1
fi
\end{lstlisting}

\section{Técnicas de Colaboración Avanzadas}

\subsection{Estrategias de Merge Avanzadas}
\begin{lstlisting}[caption=Estrategias de merge personalizadas]
# Merge con estrategia específica
git merge -X ours feature-branch     # En conflicto, preferir nuestra versión
git merge -X theirs feature-branch   # En conflicto, preferir su versión

# Merge ignorando whitespace
git merge -X ignore-space-change feature-branch

# Merge con estrategia personalizada
git merge -s recursive -X patience feature-branch

# Merge octopus (múltiples branches)
git merge branch1 branch2 branch3

# Crear merge commit sin fast-forward
git merge --no-ff feature-branch

# Merge sin commit automático
git merge --no-commit feature-branch
\end{lstlisting}

\subsection{Gestión Avanzada de Ramas}
\begin{lstlisting}[caption=Operaciones complejas con ramas]
# Ver ramas merged/no-merged
git branch --merged main
git branch --no-merged main

# Eliminar ramas remotas huérfanas
git remote prune origin

# Renombrar rama actual
git branch -m nuevo-nombre

# Copiar rama
git branch nueva-rama rama-existente

# Crear rama desde commit específico
git branch nueva-rama commit-hash

# Seguir rama remota
git branch --set-upstream-to=origin/feature-branch

# Ver información detallada de ramas
git branch -vv
\end{lstlisting}

\section{Git para Proyectos Empresariales Java}

\subsection{Monorepos y Submódulos}
\begin{lstlisting}[caption=Gestión de proyectos grandes]
# Añadir submódulo
git submodule add https://github.com/empresa/libreria.git libs/libreria

# Clonar con submódulos
git clone --recurse-submodules proyecto.git

# Actualizar submódulos
git submodule update --remote

# Inicializar submódulos después de clon
git submodule init
git submodule update

# Trabajar dentro de submódulo
cd libs/libreria
git checkout master
# hacer cambios
git commit -am "Actualizar librería"

# Actualizar referencia en proyecto principal
cd ../..
git add libs/libreria
git commit -m "Actualizar submódulo librería"
\end{lstlisting}

\subsection{Git Worktree - Múltiples Working Directories}
\begin{lstlisting}[caption=Trabajar en múltiples ramas simultáneamente]
# Crear worktree adicional
git worktree add ../hotfix-workspace hotfix/critical-bug

# Listar worktrees
git worktree list

# Trabajar en paralelo
cd ../hotfix-workspace
# hacer cambios de hotfix

# En workspace principal seguir con feature
# hacer cambios de feature

# Eliminar worktree
git worktree remove ../hotfix-workspace
\end{lstlisting}

\section{Técnicas de Debugging Avanzadas}

\subsection{Git Reflog - Historial de Referencias}
\begin{lstlisting}[caption=Recuperar commits perdidos]
# Ver historial de HEAD
git reflog

# Ver reflog de rama específica
git reflog show feature-branch

# Recuperar commit "perdido"
git checkout commit-hash-from-reflog
git branch recovered-branch

# Limpiar reflog (cuidado!)
git reflog expire --expire=30.days refs/heads/main
git gc --prune=30.days.ago
\end{lstlisting}

\subsection{Git Stash Avanzado}
\begin{lstlisting}[caption=Gestión avanzada de stash]
# Stash con mensaje
git stash push -m "WIP: implementando feature X"

# Stash incluyendo archivos no rastreados
git stash -u

# Stash archivos específicos
git stash push path/to/file.java

# Ver contenido de stash
git stash show -p stash@{0}

# Aplicar stash sin eliminarlo
git stash apply stash@{0}

# Crear rama desde stash
git stash branch nueva-rama stash@{0}

# Limpiar stash
git stash clear
\end{lstlisting}

\section{Optimización y Performance}

\subsection{Git LFS para Archivos Grandes}
\begin{lstlisting}[caption=Gestión de archivos grandes]
# Instalar Git LFS
git lfs install

# Trackear archivos por extensión
git lfs track "*.jar"
git lfs track "*.war"

# Trackear archivos por path
git lfs track "docs/assets/*"

# Ver archivos trackeados
git lfs track

# Migrar archivos existentes
git lfs migrate import --include="*.jar"

# Ver información de LFS
git lfs ls-files
git lfs status
\end{lstlisting}

\subsection{Mantenimiento del Repositorio}
\begin{lstlisting}[caption=Optimización y limpieza]
# Garbage collection manual
git gc --aggressive --prune=now

# Verificar integridad
git fsck --full

# Ver estadísticas del repositorio
git count-objects -v

# Limpiar referencias remotas
git remote prune origin

# Repack para optimizar
git repack -a -d --depth=250 --window=250

# Limpiar archivos no rastreados
git clean -fdx
\end{lstlisting}

\section{Integración con Herramientas Java}

\subsection{Maven y Git}
\begin{lstlisting}[caption=Integración con Maven]
# Plugin de release de Maven con Git
mvn release:prepare release:perform

# Configuración en pom.xml
<scm>
    <connection>scm:git:git://github.com/empresa/proyecto.git</connection>
    <developerConnection>scm:git:ssh://git@github.com/empresa/proyecto.git</developerConnection>
    <url>https://github.com/empresa/proyecto</url>
</scm>

# Git ignore específico para Maven
echo "target/" >> .gitignore
echo "*.releaseBackup" >> .gitignore
echo "release.properties" >> .gitignore
\end{lstlisting}

\subsection{Gradle y Git}
\begin{lstlisting}[caption=Integración con Gradle]
# Plugin de release para Gradle
plugins {
    id 'net.researchgate.release' version '2.8.1'
}

# Configurar versión desde Git
version = getVersionFromGit()

def getVersionFromGit() {
    def stdout = new ByteArrayOutputStream()
    exec {
        commandLine 'git', 'describe', '--tags', '--always'
        standardOutput = stdout
    }
    return stdout.toString().trim()
}
\end{lstlisting}

\section{Scripts y Automatización Avanzada}

\subsection{Scripts de Git Personalizados}
\begin{lstlisting}[caption=Aliases y scripts avanzados]
# Configurar aliases avanzados
git config --global alias.cleanup-branches '!git branch --merged main | grep -v "main\|develop" | xargs -n 1 git branch -d'

git config --global alias.find-merge '!sh -c '\''merge=$(git rev-list $1..HEAD --ancestry-path | tail -1) && git show $merge'\'' -'

git config --global alias.contributors 'shortlog -sn --all --no-merges'

git config --global alias.changelog 'log --oneline --decorate --graph --since="1 week ago"'

# Script para crear feature branch
#!/bin/bash
# git-new-feature
FEATURE_NAME=$1
if [ -z "$FEATURE_NAME" ]; then
    echo "Uso: git new-feature nombre-feature"
    exit 1
fi

git checkout develop
git pull origin develop
git checkout -b feature/$FEATURE_NAME
git push -u origin feature/$FEATURE_NAME
\end{lstlisting}

\section{Casos de Uso Avanzados Específicos}

\subsection{Migración Compleja desde SVN}
\begin{lstlisting}[caption=Migración avanzada con preservación de historia]
# Crear mapeo completo de usuarios
svn log -q | awk -F '|' '/^r/ {gsub(/ /, "", $2); sub(" $", "", $2); print $2" = "$2" <"$2"@empresa.com>"}' | sort -u > authors-transform.txt

# Migración con todas las ramas y tags
git svn clone --stdlayout --authors-file=authors-transform.txt \
    --prefix=svn/ http://svn.empresa.com/repo/proyecto

# Convertir tags SVN a tags Git reales
for t in $(git for-each-ref --format='%(refname:short)' refs/remotes/svn/tags); do
    git tag ${t/svn\/tags\//} $t
    git branch -r -d $t
done

# Limpiar referencias SVN
git for-each-ref --format='delete %(refname)' refs/remotes/svn | git update-ref --stdin
git remote rm svn 2>/dev/null
\end{lstlisting}

\subsection{Deployment con Git}
\begin{lstlisting}[caption=Estrategias de deployment]
# Git deployment con hooks
#!/bin/bash
# post-receive hook en servidor

# Checkout automático al recibir push
cd /var/www/aplicacion
git --git-dir=/var/repo/proyecto.git --work-tree=/var/www/aplicacion checkout -f

# Ejecutar build
mvn clean package -DskipTests

# Reiniciar servicio
systemctl restart aplicacion.service

echo "Deployment completado exitosamente"
\end{lstlisting}

Con esta extensión profunda, ahora tienes una comprensión completa de Git desde sus fundamentos internos hasta las técnicas más avanzadas para proyectos empresariales Java.

\end{document}
